\chapter{Surgery on a $\delta$-neck}
\label{chap:surgery-delta-neck}

\section{Notation and the Statement of the Result}\label{necksurgery}

In this chapter we describe the surgery process. For this chapter we
fix:

\begin{enumerate}
\item[(1)] A $\delta$-neck $(N,g)$ centered at a point $x_0$.
We denote by $\rho\colon  S^2\times (-\delta^{-1},\delta^{-1})\to N$
the diffeomorphism that gives the $\delta$-neck structure.
\item[(2)] Standard initial conditions $(\Ar^3,g_0)$.
\end{enumerate}
We denote by $h_0\times ds^2$ the metric on $S^2\times \Ar$ which is
the product of the round metric $h_0$ on $S^2$ of scalar curvature
$1$ and the Euclidean metric $ds^2$ on $\Ar$. We denote by
$N^-\subset N$ the image $\rho((-\delta^{-1},0]\times S^2)$ and we
denote by $s\colon N^-\to (-\delta^{-1},0]$ the composition
$\rho^{-1}$ followed by the projection to the second factor.


Recall that the standard initial metric $(\Ar^3,g_0)$ is invariant
under the standard $SO(3)$-action on $\Ar^3$. We let $p_0$ denote
the origin in $\Ar^3$. It is  the fixed point of this action and is
called the  {\em tip}
of the standard initial metric. Recall from Lemma~\ref{A_0} that
there are $A_0>0$ and an isometry
$$\psi\colon (S^2\times
(-\infty,4],h_0\times ds^2)\to (\Ar^3\setminus B(p_0,A_0),g_0).$$ The
composition of $\psi^{-1}$ followed by projection onto the second
factor defines a map $s_1\colon \Ar^3\setminus B(p_0,A_0)\to
(-\infty,4]$. Lastly, there is $0<r_0<A_0$ such that on $B(p_0,r_0)$
the metric $g_0$ is of constant sectional curvature $1/4$. We extend
the map $s_1$ to a continuous map $s_1\colon \Ar^3\to
(-\infty,4+A_0]$ defined by $s_1(x)=A_0+4-d_{g_0}(p,x)$.  This map
is an isometry along each radial geodesic ray emanating from $p_0$.
It is smooth except at $p_0$ and sends $p_0$ to $4+A_0$. The
pre-images of $s_1$ on $(-\infty,4+A_0)$ are $2$-spheres with round
metrics of scalar curvature at least $1$.




The surgery process is a local
one defined on the $\delta$-neck $(N,g)$. The surgery process
replaces  $(N,g)$ by a smooth Riemannian manifold $({\mathcal
S},\widetilde g)$. The underlying smooth manifold ${\mathcal S}$ is
obtained by gluing together $\rho(S^2\times(-\delta^{-1},4))$ and
$B(p_0,A_0+4)$ by identifying $\rho(x,s)$ with $\psi(x,s)$ for all
$x\in S^2$ and all $s\in (0,4)$. The functions $s$ on $ N^-$ and
$s_1$ agree on their overlap and hence together define a function
$s\colon {\mathcal S}\to (-\delta^{-1},4+A_0]$, a function smooth
except at $p_0$. In order to define the metric $\widetilde g$ we
must make some universal choices. We fix once and for all two bump
functions $\alpha\colon [1,2]\to [0,1]$, which is required to be
identically $1$ near $1$ and identically $0$ near $2$, and
$\beta\colon [4+A_0-r_0,4+A_0]\to [0,1]$, which is required to be
identically $1$ near $4+A_0-r_0$ and identically $0$ on
$[4+A_0-r_0/2,A_0]$. These functions are chosen once and for all and
are independent of $\delta$ and $(N,g)$. Next we set
$\eta=\sqrt{1-\delta}$.
The purpose of this choice is the following:

\begin{claim}[Distance-decreasing map from delta-neck to surgery cap model]
\label{claim:neck-to-cap-distance-decreasing}\label{distdecrease}
Let $\xi\colon N\to \Ar^3 $ be the map that sends $\rho(S^2\times
[A_0+4,\delta^{-1}))$ to the origin $0\in \Ar^3$ (i.e., to the tip of
the surgery cap) and for every $s<A_0+4$ sends $(x,s)$ to the point
in $\Ar^3$ in the radial direction $x$ from the origin at
$g_0$-distance $A_0+4-s$. Then $\xi$ is a distance decreasing map
from $(N,R(x_0)g)$ to $(\Ar^3,\eta g_0)$.
\end{claim}

\begin{proof}
\uses{lem:standard-initial-metric-distance-decreasing}
Since $R(x_0)g$ is within $\delta$ of $h_0\times ds^2$, it follows
that $R(x_0)g\ge \eta (h_0\times ds^2)$. But according to
Lemma~\ref{decrease} the map $\xi$ given in the statement of the
claim is a distance non-increasing map from $h_0\times ds^2$ to
$g_0$. The claim follows immediately.
\end{proof}


 The last choices we need to make are of
constants $C_0<\infty$ and $q<\infty$, with $C_0 \gg q$, but both of these are
independent of $\delta$. These choices will be made later. Given all these
choices, we define a function
$$f(s)=\begin{cases} 0 & s\le 0 \\ C_0\delta e^{-q/s} & s>0,\end{cases}$$
and then we define the metric $\tilde g$ on ${\mathcal S}$ by first
defining a metric:
$$\hat g=\begin{cases} \mathrm{exp}(-2f(s))R(x_0)\rho^*g & \ \ \mathrm{on} \ \
s^{-1}(-\infty,1] \\
\mathrm{exp}(-2f(s))\left(\alpha(s)R(x_0)\rho^*g+(1-\alpha(s))\eta
g_0\right)  & \ \ \mathrm{on}\ \  s^{-1}([1,2]) \\
\mathrm{exp}(-2f(s))\eta g_0 & \ \ \mathrm{on}\ \ s^{-1}([2,A_{r_0}]\\
\left[\beta(s)\mathrm{exp}(-2f(s))+(1-\beta(s))\mathrm{exp}(-2f(4+A_0))\right]\eta
g_0 & \ \ \mathrm{on} \ \ s^{-1}([A_{r_0},A']),\end{cases}$$ where
$A_{r_0}=4+A_0-r_0$ and $A'=A_0+4$.
 Then we define
$$\widetilde g=R(x_0)^{-1}\hat g.$$








\begin{theorem}[Local surgery on a delta-neck: curvature control after surgery]
\label{thm:local-surgery}
\uses{claim:neck-to-cap-distance-decreasing}\label{LOCALSURGERY}
There are constants $C_0,q,R_0<\infty$ and $\delta'_0>0$ such that the
following hold for the result $({\mathcal S},\widetilde g)$ of surgery on
$(N,g)$ provided that $R(x_0)\ge R_0$, $0<\delta\le \delta'_0$. Define $f(s)$
as above with the constants $C_0,\delta$ and then use $f$ to define surgery on
a $\delta$-neck $N$ to produce $({\mathcal S},\tilde g)$. Then the following
hold.
\begin{itemize}
\item Fix $t\ge 0$. For any $p\in N$,
let $X(p)=\mathrm{max}(0,-\nu_{g}(p))$, where $\nu_ g(p)$ is the
smallest eigenvalue of $Rm_{ g}(p)$. Suppose that for all $p\in N$
we have:
\begin{enumerate}\item[(1)] $R(p)\ge \frac{-6}{1+4t}$, and
\item[(2)] $R(p)\geq 2X(p)\left(\mathrm{log}X(p)+\mathrm{log}(1+t)-3\right)$, whenever $0<X(p)$.
\end{enumerate}
Then the curvature of $({\mathcal S},\widetilde g)$ satisfies the
same equations at every point of ${\mathcal S}$ with the same value
of $t$.
\item The restriction of the metric $\widetilde g$ to $s^{-1}([1,4+A_0])$ has
positive sectional curvature.
\item Let $\xi\colon N\to {\mathcal S}$
be the map given in Claim~\ref{distdecrease}. Then it is a distance
decreasing map from $g$ to $\tilde g$.
\item For any $\delta''>0$ there is $\delta'_1=\delta'_1(\delta'')>0$  such that if
$\delta\le \mathrm{min}(\delta'_1,\delta'_0)$, then  the restriction of
$\hat g$ to $B_{\hat g}(p_0,(\delta'')^{-1})$ in $({\mathcal S},\hat
g)$ is $\delta''$-close in the $C^{[1/\delta'']}$-topology to the
restriction of the standard initial metric $g_0$ to
$B_{g_0}(p_0,(\delta'')^{-1})$.
\end{itemize}
\end{theorem}




The rest of this chapter is devoted to the proof of this theorem.




Before starting the curvature computations let us make a remark about the surgery cap.

\begin{definition}[Surgery cap]
\label{def:surgery-cap}\label{surgerycap}
The image in ${\mathcal S}$ of $B_{g_0}(p_0,0,A_0+4)$ is called the {\em
surgery cap}.
\end{definition}

The following is immediate from the definitions provided that
$\delta>0$ is sufficiently small.

\begin{lemma}[Surgery cap lies between two concentric metric balls]
\label{lem:surgery-cap-ball-bounds}
\uses{def:surgery-cap}\label{capball}
The surgery cap in $({\mathcal S},\tilde g)$ has a metric that
differs from the one coming from a rescaled version of the standard
solution. Thus, the image of this cap is not necessarily a metric
ball. Nevertheless for $\epsilon<1/200$ the image of this cap will
be contained in the metric ball in ${\mathcal S}$ centered at $p_0$
of radius $R(x_0)^{-1/2}(A_0+5)$ and will contain the metric ball
centered at $p_0$ of radius $R(x_0)^{-1/2}(A_0+3)$. Notice also that
the complement of the closure of the surgery cap in ${\mathcal S}$
is isometrically identified with $N^-$.
\end{lemma}


\section{Preliminary computations}

We shall compute  in a slightly more general setup.
 Let $I$ be an open interval contained in $(-\delta^{-1},4+A_0)$ and let
$h$ be a metric  on $S^2\times I$ within $\delta$ in the
$C^{[1/\delta]}$-topology of the restriction to this open
submanifold of the standard metric $h_0\times ds^2$.  We let $\hat
h=e^{-2f}h$. Fix local coordinates near a point $y\in S^2\times I$.
We denote by $\nabla$ the covariant derivative for $h$ and by
$\widehat\nabla$ the covariant derivative for $\hat h$.
 We also denote by
$(R_{ijkl})$ the matrix of the Riemann curvature operator of $h$
in the associated basis of $\wedge^2T(S^2\times I)$ and by $(\hat
R_{ijkl})$ the matrix of the Riemann curvature operator of $\hat
h$ with respect to the same basis. Recall the formula for the
curvature of a conformal change of metric (see, (3.34) on p.51 of
\cite{Sakai}):
\begin{eqnarray}\label{confchange}\hat R_{ijkl} & = &
e^{-2f}\left(R_{ijkl}-f_jf_kh_{il}+f_jf_lh_{ik}+f_if_kh_{jl}-f_if_lh_{jk}\right.
\\ & & \left.-
(\wedge^2h)_{ijkl}|\nabla
f|^2-f_{jk}h_{il}+f_{ik}h_{jl}+f_{jl}h_{ik}-f_{il}h_{jk}\right).\nonumber\end{eqnarray}
Here, $f_i$ means $\partial_if$,
$$f_{ij}=\Hess_{ij}(f)=\partial_if_j-f_l\Gamma_{ij}^l,$$
and $\wedge^2h$ is the metric induced by $h$ on $\wedge^2TN$, so
that
$$\wedge^2h_{ijkl}=h_{ik}h_{jl}-h_{il}h_{jk}.$$

Now we introduce the notation $O(\delta)$. When we say that a
quantity is $O(\delta)$ we mean that there is some universal
constant $C$ such that, provided that $\delta>0$ is sufficiently
small, the absolute value of the quantity is $\le C\delta$. The
universal constant is allowed to change from inequality to
inequality.


In our case we take local coordinates adapted to the $\delta$-neck:
$(x^0,x^1,x^2)$ where $x^0$ agrees with the $s$-coordinate and
$(x^1,x^2)$ are Gaussian local coordinates on the $S^2$ such that
$dx^1$ and $dx^2$ are orthonormal at the point in question in the
round metric $h_0$. The function $f$ is a function only of $x^0$.
Hence $f_i=0$ for $i=1,2$. Also, $f_0=\frac{q}{s^2}f$. It follows
that
$$|\nabla f|_h=\frac{q}{s^2}f\cdot(1+O(\delta)),$$ so that
$$|\nabla
f|_h^2=\frac{q^2}{s^4}f^2\cdot (1+O(\delta)).$$

Because the metric $h$ is $\delta$-close to the product $h_0\times
ds^2$, we see that $h_{ij}(y)=(h_0)_{ij}(y)+O(\delta)$ and the
Christoffel symbols $\Gamma_{ij}^k(y)$ of $h$ are within $\delta$ in
the $C^{([1/\delta]-1)}$-topology of those of the product metric
$h_0\times ds^2$. In particular, $\Gamma_{ij}^0=O(\delta)$ for all
$ij$.
  The components $f_{ij}$ of the Hessian with
respect to $h$ are given by
$$f_{00}=\left(\frac{q^2}{s^4}-\frac{2q}{s^3}\right)f+\frac{q}{s^2}fO(\delta),$$
$$f_{i0}=\frac{q}{s^2}fO(\delta)\ \ \ \mathrm{for}\ \ 1\le i\le 2,$$
$$f_{ij}=\frac{q}{s^2}fO(\delta)\ \ \ \mathrm{for}\ \ 1\le i,j\le 2.$$
In the following $a,b,c,d$ are indices taking values $1$ and $2$.
Substituting in Equation~(\ref{confchange}) yields
\begin{eqnarray*}\hat
R_{0a0b} & = &
e^{-2f}\left(R_{0a0b}+\frac{q^2}{s^4}f^2h_{ab}-h_{ab}(\frac{q^2}{s^4})f^2(1+O(\delta))
+\left(\frac{q^2}{s^4}-\frac{2q}{s^3}\right)fh_{ab} \right. \\
& & \left. +\frac{q}{s^2}fO(\delta)\right)\\
& = &
e^{-2f}\left(R_{0a0b}+\left(\frac{q^2}{s^4}-\frac{2q}{s^3}\right)fh_{ab}+
h_{ab}(\frac{q^2}{s^4})f^2O(\delta)
+\frac{q}{s^2}fO(\delta)\right)\end{eqnarray*} Also, we have
\begin{eqnarray*}\hat
R_{ab0c} & = &
e^{-2f}\left(R_{ab0c}-(\wedge^2h)_{ab0c}(\frac{q^2}{s^4})f^2(1+O(\delta))
+\frac{q}{s^2}fO(\delta)\right) \\
& = & e^{-2f}\left(R_{ab0c}+(\frac{q^2}{s^4})f^2O(\delta)
+\frac{q}{s^2}fO(\delta)\right).\end{eqnarray*}
 Lastly,
\begin{eqnarray*} \hat R_{1212} & = &
e^{-2f}\left(R_{1212}-(\wedge^2h)_{1212}\frac{q^2}{s^4}f^2(1+O(\delta))+\frac{q}{s^2}fO(\delta)\right)
\\
& = &
e^{-2f}\left(R_{1212}-\frac{q^2}{s^4}f^2(1+O(\delta))+\frac{q}{s^2}fO(\delta)\right).
\end{eqnarray*}


Now we are ready to fix the constant $q$. We fix it so that for all $s\in
[0,4+A_0]$ we have \begin{equation}\label{qeqn}
 q \gg (4+A_0)^2\ \ \ \mathrm{and}\ \ \ \frac{q^2}{s^4}e^{-q/s} \ll 1.
 \end{equation}
It follows immediately that $q^2/s^4\gg q/s^3$ for all $s\in [0,4+A_0]$.
 We are not yet ready
to fix the constant $C_0$, but once we do we shall always require $\delta$ to
satisfy
$\delta\ll C_0^{-1}$ so that for all $s\in [0,4+A_0]$ we have
$$\frac{q}{s^2}f^2\ll \frac{q^2}{s^4}f^2\ll \frac{q}{s^2}f\ll 1.$$
(These requirements are not circular, since $C_0$ and $q$ are chosen
independent of $\delta$.)




  Using
these inequalities and putting our computations in matrix form show
the following.

\begin{corollary}[Curvature formula for the conformally rescaled neck metric]
\label{cor:conformal-curvature-formula}\label{delta0}
There is $\delta'_2>0$, depending on $C_0$ and $q$,
such that if $\delta\le \delta'_2$ then we have
\begin{equation}\label{maineqn}
\left(\hat
R_{ijkl}\right)=e^{-2f}\left[\left(R_{ijkl}\right)+\begin{pmatrix}
-\frac{q^2}{s^4}f^2 & 0  \\ 0 &
\left(\frac{q^2}{s^4}-\frac{2q}{s^3}\right)f\begin{pmatrix} 1 & 0
\\ 0 & 1\end{pmatrix}
\end{pmatrix}+\left(\frac{q}{s^2}fO(\delta)\right)\right].\end{equation}
\end{corollary}

Similarly, we have the equation relating scalar curvatures
$$\hat R=e^{2f}\left(R+4\triangle f-2|\nabla f|^2\right),$$
and hence
$$\hat R=e^{2f}\left(R+4\left(\frac{q^2}{s^4}-\frac{2q}{s^3}\right)f-2\frac{q^2}{s^4}f^2+
\frac{q}{s^2}fO(\delta)\right).$$

\begin{corollary}[Scalar curvature does not decrease under the surgery conformal factor]
\label{cor:scalar-curvature-nondecreasing-conformal}
\uses{cor:conformal-curvature-formula}\label{Rcompar}
For any constant $C_0<\infty$ and any $\delta<\mathrm{min}(\delta_2',C_0^{-1})$ we
have $\hat R\ge R$.
\end{corollary}

\begin{proof}
By our choice of $q$, since $C_0\delta<1$, then $f^2\ll f$ and $q^2/s^4\gg \mathrm{max}(q/s^3,q/s^2)$ so that the result follows immediately from the above
formula.
\end{proof}


Now let us compute the eigenvalues of the curvature $ R_{ijkl}(y)$
for any $y\in S^2\times I$.

\begin{lemma}[Near-diagonal curvature eigenbasis for a near-cylindrical metric]
\label{lem:near-cylinder-curvature-eigenbasis}\label{newbasis}
There is a $\delta_3'>0$ such that the following hold if $\delta\le
\delta_3'$. Let $\{e_0,e_1,e_2\}$ be an orthonormal basis for the
tangent space at a point $y\in S^2\times I$ for the metric
$h_0\times ds^2$ with the property that $e_0$ points in the
$I$-direction. Then there is a basis $\{f_0,f_1,f_2\}$ for this
tangent space so that the following hold:
\begin{enumerate}
\item[(1)] The basis is orthonormal in the metric $h$.
\item[(2)] The change of basis matrix expressing the $\{f_0,f_1,f_2\}$
in terms of $\{e_0,e_1,e_2\}$ is of the form $\Id+O(\delta)$.
\item[(3)] The Riemann curvature of $h$ in the basis $\{f_0\wedge f_1,f_1\wedge f_2,f_2\wedge f_0\}$
of $\wedge^2T_y(S^2\times I)$ is
$$\begin{pmatrix} 1/2 & 0 & 0 \\ 0 & 0 & 0 \\ 0 & 0 & 0
\end{pmatrix}+O(\delta).$$
\end{enumerate}
\end{lemma}


\begin{proof}
Since $h$ is within $\delta$ of $h_0\times ds^2$ in the
$C^{[1/\delta]}$-topology, it follows that the matrix for $h(y)$ in
$\{e_0,e_1,e_2\}$ is within $O(\delta)$ of the identity matrix, and the matrix
for the curvature of $h$ in the associated basis of $\wedge^2T_y(S^2\times I)$
is within $O(\delta)$ of the curvature matrix for $h_0\times ds^2$, the latter
being the diagonal matrix with diagonal entries $\{1/2,0,0\}$. Thus, the usual
Gram-Schmidt orthonormalization process constructs the basis $\{f_0,f_1,f_2\}$
satisfying the first two items. Let $A=(A^{ab})$ be the change of basis matrix
expressing the $\{f_a\}$ in terms of the $\{e_b\}$, so that $A=\Id+O(\delta)$. The curvature of $h$ in this basis is then given by
$B^{\tr}(R_{ijkl})B$ where $B=\wedge^2A$ is the induced change of basis matrix
expressing the basis $\{f_0\wedge f_1,f_1\wedge f_2,f_2\wedge f_0\}$ in terms
of $\{e_0\wedge e_1,e_1\wedge e_2,e_2\wedge e_0\}$. Hence, in the basis
$\{f_0\wedge f_1,f_1\wedge f_2,,f_2\wedge f_0\}$ the curvature matrix for $h$
is within $O(\delta)$ of the same diagonal matrix. For $\delta$ sufficiently
small then the eigenvalues of the curvature matrix for $h$ are within
$O(\delta)$ of $(1/2,0,0)$.
\end{proof}

\begin{corollary}[Curvature block-diagonalizes in the adapted orthonormal basis]
\label{cor:curvature-diagonal-block-basis}
\uses{lem:near-cylinder-curvature-eigenbasis}\label{2ndnewbasis}
The following holds provided that $\delta\le \delta_3'$. It is
possible to choose the basis $\{f_0,f_1,f_2\}$ satisfying the
conclusions of Lemma~\ref{newbasis} so that in addition the
curvature matrix for $(R_{ijkl}(y))$ is of the form
$$\begin{pmatrix} \lambda & 0 & 0 \\ 0 & \alpha & \beta \\ 0 & \beta & \gamma\end{pmatrix}$$
with $|\lambda-\frac{1}{2}|\le O(\delta)$ and
$|\alpha|,|\beta|,|\gamma|\le O(\delta)$.
\end{corollary}


\begin{proof}
We have an $h$-orthonormal basis $\{f_0\wedge f_1,f_1\wedge
f_2,f_2\wedge f_0\}$ for $\wedge^2T_y(S^2\times \Ar)$ in which the
quadratic form $(R_{ijkl}(y)$ is
$$\begin{pmatrix} 1/2 & 0 & 0 \\ 0 & 0 & 0 \\ 0 & 0 & 0\end{pmatrix}
+O(\delta).$$ It follows that the restriction to the $h$-unit sphere
in $\wedge^2T_y(S^2\times \Ar)$ of this quadratic form achieves its
maximum value at some vector $v$, which, when written out in this
basis, is given by $(x,y,z)$ with $|y|,|z|\le O(\delta)$ and
$|x-1|\le O(\delta)$. Of course, this maximum value is within
$O(\delta)$ of $1/2$. Clearly, on the $h$-orthogonal subspace to
$v$, the quadratic form is given by a matrix all of whose entries
are $O(\delta)$ in absolute value. This gives us a new basis of
$\wedge^2T_y(S^2\times I)$ within $O(\delta)$ of the given basis in
which $(R_{ijkl}(y))$ is diagonal. The corresponding basis for
$T_y(S^2\times\Ar)$ is as required.
\end{proof}

Now we consider the expression  $(\hat R_{ijkl}(y))$ in this basis.

\begin{lemma}[Curvature matrix of the conformally rescaled metric in the adapted basis]
\label{lem:rhat-matrix-adapted-basis}
\uses{cor:curvature-diagonal-block-basis}\label{Rhatmatrix} Set $\delta'_4=\mathrm{min}(\delta_2',\delta_3')$.  Suppose
that $\delta\le \mathrm{min}(\delta'_4,C_0^{-1})$. Then in the basis
$\{f_0,f_1,f_2\}$ for $T_y(S^2\times I)$ as in Corollary~\ref{2ndnewbasis} we
have
$$(\hat R_{ijkl}(y)) =e^{-2f}\left[\begin{pmatrix} \lambda & 0 & 0
\\ 0 & \alpha & \beta \\ 0 & \beta  & \gamma\end{pmatrix}
+\begin{pmatrix} -\frac{q^2}{s^4}f^2 & 0  \\ 0 &
\left(\frac{q^2}{s^4}-\frac{2q}{s^3}\right)f\begin{pmatrix} 1 & 0
\\ 0 & 1\end{pmatrix}
\end{pmatrix}+\left(\frac{q^2}{s^4}fO(\delta)\right)\right]$$
where $\lambda,\alpha,\beta,\gamma$ are the constants in
Lemma~\ref{2ndnewbasis} and the first matrix is the expression for
$(R_{ijkl}(y))$ in this basis.
\end{lemma}


\begin{proof}
\uses{cor:conformal-curvature-formula}
We simply conjugate the expression in Equation~(\ref{maineqn}) by the change of
basis matrix and use the fact that by our choice of $q$ and the fact that
$C_0\delta<1$, we have $f\gg  f^2$ and $q/s^3\ll q^2/s^4$.
\end{proof}

\begin{corollary}[Fully diagonalized curvature matrices for the neck metric and its rescaling]
\label{cor:diagonalized-curvature-matrices}\label{13.10}
\uses{lem:near-cylinder-curvature-eigenbasis,lem:rhat-matrix-adapted-basis}
Assuming that $\delta\le \mathrm{min}(\delta'_4,C_0^{-1})$, there is an
$h$-orthonormal basis $\{f_0,f_1,f_2\}$ so that in the associated basis for
$\wedge^2T_y(S^2\times I)$ the matrix $(R_{ijkl}(y))$ is diagonal and given by
$$\begin{pmatrix} \lambda & 0 & 0 \\ 0 & \mu & 0 \\ 0 & 0 &
\nu\end{pmatrix}$$ with $|\lambda-1/2|\le O(\delta)$ and
$|\mu|,|\nu|\le O(\delta)$. Furthermore, in this same basis the
matrix  $(\hat R_{ijkl}(y))$ is
$$e^{-2f}\left[\begin{pmatrix}\lambda & 0 & 0 \\ 0 & \mu & 0 \\ 0 & 0 &
\nu\end{pmatrix}+\begin{pmatrix} -\frac{q^2}{s^4}f^2 & 0  \\ 0 &
\left(\frac{q^2}{s^4}-\frac{2q}{s^3}\right)f\begin{pmatrix} 1 & 0 \\
0 & 1\end{pmatrix}\end{pmatrix}+\frac{q^2}{s^4}fO(\delta)\right].$$
\end{corollary}

\begin{proof}
\uses{lem:near-cylinder-curvature-eigenbasis}
To diagonalize $(R_{ijkl}(y))$ we need only rotate in the
$\{f_1\wedge f_2,f_2\wedge f_3\}$-plane. Applying this rotation to
the expression in Lemma~\ref{newbasis} gives the result.
\end{proof}


\begin{corollary}[Eigenvalue estimates for the conformally rescaled curvature]
\label{cor:conformal-eigenvalue-estimates}\label{improve}
There is a constant $A<\infty$ such that the following holds for the given value of $q$
and any $C_0$ provided that $\delta$ is sufficiently small.
Suppose that the eigenvalues for the curvature matrix of $h$ at $y$
are $\lambda\ge \mu\ge \nu$.
Then the eigenvalues for the curvature of $\hat h$ at the point $y$
are given by $\lambda',\mu',\nu'$, where
$$\left|\lambda'-e^{2f}\left(\lambda-\frac{q^2}{s^4}f^2\right)\right|
\le \frac{q^2}{s^4}fA\delta$$
$$\left|\mu'-e^{2f}\left(\mu+\left(\frac{q^2}{s^4}-\frac{2q}{s^3}\right)f\right)\right|\le
\frac{q^2}{s^4}fA\delta$$
$$\left|\nu'-e^{2f}\left(\nu+\left(\frac{q^2}{s^4}-\frac{2q}{s^3}\right)f\right)\right|\le
\frac{q^2}{s^4}fA\delta.$$

In particular, we have
$$\nu'\ge e^{2f}\left(\nu+\frac{q^2}{2s^4}f\right)$$
$$\mu'\ge e^{2f}\left(\mu+\frac{q^2}{2s^4}f\right).$$
\end{corollary}

\begin{proof}
\uses{cor:diagonalized-curvature-matrices,lem:rhat-matrix-adapted-basis}
Let $\{f_0,f_1,f_2\}$ be the $h$-orthonormal basis given in
Corollary~\ref{13.10}. Then $\{e^ff_0,e^ff_1,e^ff_2\}$ is orthonormal for $\hat
h=e^{-2f}h$. This change multiplies the curvature matrix by $e^{4f}$. Since
$f\ll 1$, $e^{4f}<2$ so that the expression for $(\hat R_{ijkl}(y))$ in this
basis is exactly the same as in Lemma~\ref{Rhatmatrix} except that the factor
in front is $e^{2f}$ instead of $e^{-2f}$. Now, it is easy to see that since
$((q^2/s^4)fA\delta)^2\ll (q^2/s^4)fA\delta$ the eigenvalues will differ from
the diagonal entries by at most a constant multiple of $(q^2/s^4)fA\delta$.


 The first three inequalities are  immediate
from the previous corollary. The last two follow since $q^2/s^4\gg q/s^3$ and
$\delta\ll 1$.
\end{proof}




One important consequence of this computation is the following:





\begin{corollary}[Smallest curvature eigenvalue increases under the surgery conformal change]
\label{cor:smallest-eigenvalue-increases}
\uses{cor:conformal-eigenvalue-estimates}\label{smeigen}
For the given value of $q$ and for any $C_0$,
assuming that $\delta>0$ is sufficiently small, the smallest
eigenvalue of $\Rm_{\hat h}$ is greater than the smallest eigenvalue
of $\Rm_h$ at the same point. Consequently, at any point where $h$
has non-negative curvature  so does $\hat h$.
\end{corollary}

\begin{proof}
Since $|\lambda-1/2|, |\mu|, |\nu|$ are all $O(\delta)$ and since
$\frac{q^2}{s^4}f\ll 1$, it follows that the smallest eigenvalue of $(\hat
R_{ijkl}(y))$ is either $\mu'$ or $\nu'$. But it is immediate from the above
expressions  that $\mu'>\mu$ and $\nu'>\nu$. This completes the proof.
\end{proof}


Now we are ready to fix $C_0$.
There is a universal constant $K$ such that for all $\delta>0$ sufficiently small
and for any $\delta$-neck $(N,h)$ of
scale one, every eigenvalue of $\Rm_h$ is at least $-K\delta$.
We set
$$C_0=2Ke^{q}.$$

\begin{lemma}[Positivity of curvature eigenvalues on the surgery region]
\label{lem:positive-eigenvalues-surgery-region}
\uses{cor:conformal-eigenvalue-estimates}\label{nuprime>0}
With these choices of $q$ and $C_0$ for any $\delta>0$ sufficiently small we have
$\nu'>0$ and $\mu'>0$ for $s\in [1,4+A_0]$ and $\lambda'>1/4$.
\end{lemma}


\begin{proof}
Then by the previous result we have
$$\nu'\ge e^{2f}\left(\nu+\frac{q^2}{2s^4}f\right).$$
It is easy to see that since $q\gg (4+A_0)$ the function $(q^2/2s^4)f$ is an
increasing function on $[1,4+A_0]$. Its value at $s=1$ is
$(q^2/2)e^{-q}C_0\delta>K\delta$. Hence $\nu+\frac{q^2}{2s^4}f>0$ for all $s\in
[1,4+A_0]$ and consequently $\nu'>0$ on this submanifold. The same argument
shows $\mu'>0$. Since $q^2/s^4 f^2\ll 1$ and $0<f$, the statement about
$\lambda'$ is immediate.
\end{proof}





\section{The proof of Theorem~{\ref{LOCALSURGERY}}}




\subsection{Proof of the first two items for $s< 4$}

We consider the metric in the region $s^{-1}(-\delta^{-1},4)$ given by
$$h=\alpha(s)R_g(x_0)\rho^*g+(1-\alpha(s))\eta g_0.$$
There is a constant $K'<\infty$ (depending on the
$C^{[1/\delta]}$-norm of  $\alpha$) such that $h$ is within
$K'\delta$ of the product metric in the
$C^{[1/(K'\delta)]}$-topology. Thus, if $\delta$ is sufficiently
small, all of the preceding computations hold with the error term
$(q^2/s^4)fAK'\delta$. Thus, provided that $\delta$ is sufficiently
small,  the conclusions about the eigenvalues hold for $e^{-2f}h$ in
the region $s^{-1}(-\delta^{-1},4)$. But $e^{-2f}h$ is exactly equal
to $R(x_0)\tilde g$ in this region. Rescaling, we conclude that on
$s^{-1}(-\delta^{-1},4)$ the smallest eigenvalue of $\tilde g$ is
greater than the smallest eigenvalue of $g$ at the corresponding
point and that $R_{\tilde g}\ge R_{g}$ in this same region.


The first conclusion of Theorem~\ref{LOCALSURGERY}  follows by
applying the above considerations to the case of
$h=R_g(x_0)\rho^*g$. Namely, we have:

\begin{proposition}[Curvature pinching is preserved by surgery on the region s<4]
\label{prop:pinching-preserved-s-less-4}
Fix $\delta>0$ sufficiently small. Suppose that for some $t\ge 0$
and every point $p\in N$ the curvature of $h$ satisfies:
\begin{enumerate}
\item[(1)] $R(p)\ge \frac{-6}{1+4t}$, and
\item[(2)] $R(p)\geq 2X(p)\left(\mathrm{log}X(p)+\mathrm{log}(1+t)-3\right)$ whenever $0<X(x,t)$.
\end{enumerate}

Then  the curvature $({\mathcal S},\widetilde g)$ satisfies the same
equation with the same value of $t$ in the region
$s^{-1}(-\delta^{-1},4)$. Also, the curvature of $\tilde g$ is
positive in the region $s^{-1}[1,4)$.
\end{proposition}

\begin{proof}
\uses{thm:local-surgery,cor:scalar-curvature-nondecreasing-conformal,lem:positive-eigenvalues-surgery-region,cor:smallest-eigenvalue-increases}
According to Corollary~\ref{smeigen}, the smallest eigenvalue of $\hat
h$ at any point $p$ is greater than or equal to the smallest
eigenvalue of $h$ at the corresponding point. According to
Corollary~\ref{Rcompar}, $\hat R(p)\ge R(p)$ for every $p\in
{\mathcal S}$. Hence, $X_{\hat h}(p)\le X_h(p)$. If $X_h(p)\ge
e^3/(1+t)$, then we have
$$\hat R(p)\ge R(p)\ge 2X_h(p)(\mathrm{log}(X_h(p)+\mathrm{log}(1+t)-3)\ge 2X_{\hat h}(p)
(\mathrm{log}(X_{\hat h}(p)+\mathrm{log}(1+t)-3).$$
If $X_h(p)<e^3(1+t)$, then $X_{\hat h}(p)<e^3/(1+t)$. Thus, in this case
since we are in a $\delta$-neck, provided that $\delta$ is sufficiently small, we have
 $R(p)\ge 0$ and hence
$$\hat R(p)\ge R(p)\ge 0>2X_{\hat p}(p)(\mathrm{log}(X_{\hat h}(p)+\mathrm{log}(1+t)-3).$$

This completes the proof in both cases.


This establishes the first item in the conclusion of
Theorem~\ref{LOCALSURGERY} for $\delta>0$ sufficiently small on
$s^{-1}(-\delta^{-1},4)$. As we have seen in Lemma~\ref{nuprime>0},
the curvature is positive on $s^{-1}[1,4)$.
\end{proof}


\subsection{Proof of the first two items for $s\ge 4$}

Now let us show that the curvature on $\tilde g$ is positive in the region
$s^{-1}([4,4+A_0])$. First of all in the preimage of the
interval $[4,4+A_0,-r_0]$ this follows
from Corollary~\ref{smeigen} and the fact that $\eta g_0$ has non-negative
curvature. As for the region $s^{-1}([4+A_0-r_0,4+A_0])$, as $\delta$ tends to
zero, the metric here tends smoothly to the restriction of the metric $g_0$ to
that subset. The metric $g_0$ has positive curvature on
$s^{-1}([4+A_0-r_0,4+A_0])$. Thus, for all $\delta>0$ sufficiently small the
metric $\tilde g$ has positive curvature on all of $s^{-1}([4+A_0-r_0,4+A_0])$.
This completes the proof of the first two items.




\subsection{Proof of the third item}

 By
construction the restriction of the metric $\widetilde g$ to
$s^{-1}((-\delta^{-1},0])$ is equal to the metric $\rho^*g$. Hence,
in this region the mapping  is an isometry. In the region
$s^{-1}([0,4])$ we have $R(x_0)\rho^*g\ge \eta g_0$ so that by
construction in this region $\rho^*g\ge \tilde g$. Lastly, in the
region $s^{-1}([4,A_0+4])$ we have $R(x_0)^{-1}\eta g_0\ge \tilde
g$. On the other hand, it follows from Lemma~\ref{distdecrease} that
the map from $([0,\delta^{-1}]\times S^2,R(x_0)\rho^* g)$ to
$(B(p_0,4+A_0),\eta g)$ is distance decreasing. This completes the
proof of the third item.

\subsection{Completion of the proof}

As $\delta$ goes to zero,  $f$ tends to zero in the
$C^\infty$-topology and $\eta$ limits to $1$. From this  the fourth
item is clear.

This completes the proof of Theorem~\ref{LOCALSURGERY}.


\section{Other properties of the result of surgery}

\begin{lemma}[Metric balls in the surgery result compared to the neck coordinate]
\label{lem:surgery-result-metric-ball-bounds}\label{stdballs}
Provided that $\delta>0$ is sufficiently small the following holds. Let $(N,g)$
be a $\delta$-neck and let $({\mathcal S},\tilde g)$ be the result of surgery
along the cental $2$-sphere of this neck. Then for any $0<D<\infty$ the ball
$B_{\tilde g}(p,D+5+A_0)\subset {\mathcal S}$ has boundary contained in
$s_N^{-1}(-(2D+2),-D/2)$.
\end{lemma}


\begin{proof}
The Riemannian manifold $({\mathcal S},\tilde g)$ is identified by a
diffeomorphism with the union of $s_{N}^{-1}(-\delta^{-1},0]$ to
$B_{g_0}(p_0,A_0+4)$ glued along their boundaries. Thus, we have a natural
identification of ${\mathcal S}$ with the ball $B_{g_0}(p,A_0+4+\delta^{-1})$
in the standard solution. This identification pulls back the metric $\tilde g$
to be within $2\delta$ of the standard initial metric. The result then follows
immediately for $\delta$ sufficiently small.
\end{proof}

